(1) $\sigma^2 = v^*vv^*v + vv^*vv^* - v^*v\,vv^* - vv^*\,v^*v = v^*v + vv^* - 0 - 0 = 1$, using $vv^*v = v$, $v^*vv^* = v^*$ and the vanishing squares. In $\Cl(1,1)$, by \claimref{std-algebra.vector-squares}: $v^*v = \tfrac14(e_0 - e_1)(e_0 + e_1) = \tfrac14(e_0^2 - e_1^2 + e_0e_1 - e_1e_0) = \tfrac14(1 + 1 + 2e_0e_1) = \tfrac12(1 + e_0e_1)$, so $\sigma = v^*v - vv^* = e_0e_1$ in this orientation; relabelling the pair reverses it. (2) From $vv^*v = v$ and the vanishing squares, $\sigma v = (v^*v - vv^*)v = -vv^*v = -v$ and $v\sigma = v(v^*v - vv^*) = vv^*v = +v$. The first relation gives $e^{\lambda\sigma/2}\,v = v\,e^{-\lambda\sigma/2}$ term by term in the exponential series, so $R_\lambda\, v\, R_\lambda^{-1} = v\,e^{-\lambda\sigma}$. Since $\sigma^2 = 1$, $e^{-\lambda\sigma} = \cosh\lambda - \sigma\sinh\lambda$ (\claimref{std-algebra.bivector-exponential}), and the second relation collapses it on the left factor: $v\,e^{-\lambda\sigma} = v(\cosh\lambda - \sinh\lambda) = e^{-\lambda}v$. The $v^*$ computation is the mirror image, with both signs reversed. (3) $\tfrac12(1 + \sigma) = \tfrac12(1 + v^*v - vv^*) = \tfrac12(v^*v + vv^* + v^*v - vv^*) = v^*v$, and likewise for the complement. (4) In the two-dimensional Fock module ($v^* = $ creation, $v = $ annihilation, $N = v^*v$): $\tfrac12\sigma = \tfrac12(v^*v - vv^*) = \tfrac12(N - (1 - N)) = N - \tfrac12$. (5) The bivector space of $\Cl(1,1)$ is the line $\R\,e_0e_1$, so every one-parameter rotor group is of the form $\lambda \mapsto \exp(c\lambda\sigma/2)$ for some real $c$: the same orbits at a rescaled parameter, which is the reparametrisation the statement reserves. The verification script checks every clause in exact arithmetic, including the finite conjugation of clause (2) at a rational rapidity.
